\subsection{Income-statement analysis}

These entries cover how non-recurring items and accounting changes are reported and adjusted, how earnings per share is computed, and how a common-size income statement and its ratios are used to evaluate performance. When revenue and expenses are recognised, and how capitalisation shifts expense between periods and cash-flow sections, are in \hyperref[meth:revenue-recognition-analysis]{revenue recognition} and \hyperref[meth:expense-recognition-capitalisation]{expense recognition}.

\begin{qfmethod}[Non-recurring items and accounting changes: material and unusual items, discontinued operations, changes in policies and estimates, error corrections, and adjusting earnings for analysis]\label{meth:non-recurring-items}
\qffield{Statement} \emph{Presentation.} IFRS prohibits presenting any income or expense as extraordinary \citep[IAS~1 para.~87]{eu-2023-1803-ias1}. Material items are disclosed separately by nature and amount (para.~97). Examples include write-downs of inventory or PP\&E and their reversals, restructurings, disposals of PP\&E and investments, discontinued operations, litigation settlements and reversals of provisions (para.~98). US GAAP eliminated the extraordinary-items concept for fiscal years beginning after 15~December 2015, retaining presentation guidance for items that are unusual or infrequent \citep[p.~2]{us-extraordinary-items-2015}.

A \emph{discontinued operation} is a component that has been disposed of or is held for sale and that is a separate major line of business or geographical area, part of a single plan to dispose of one, or a subsidiary acquired exclusively for resale \citep[IFRS~5 para.~32]{eu-2023-1803-ifrs5}. Its post-tax result and disposal gain or loss are shown as a single amount, analysed in the notes or on the face (para.~33).

\emph{Accounting changes} \citep[IAS~8]{eu-2023-1803-ias8}:
\begin{itemize}
\item a policy changes only if a standard requires it or the change gives reliable and more relevant information (para.~14); it applies retrospectively unless a transition provision says otherwise, adjusting the opening equity of the earliest period presented (paras.~19 and~22);
\item an estimate is a monetary amount subject to measurement uncertainty (para.~5); a change in estimate is recognised prospectively (para.~36) and its nature and amount disclosed (para.~39);
\item a material prior-period error is corrected retrospectively by restating comparatives (para.~42), with the per-line and EPS effects disclosed (para.~49).
\end{itemize}

\emph{Analysis.} Recast the income statement to separate recurring from transitory components, tax-effecting each item so the total still reconciles to profit; then assign items to the periods they belong to \citep[pp.~602--603]{financial-statement-analysis-2009}. Transitory items are excluded from \emph{permanent} (sustainable) income but included in \emph{economic} income, and they reveal risk exposures that may recur (p.~338). Remove a discontinued operation's results from current and past income, and its assets and liabilities from the balance sheet (p.~340).

\qffield{Notation} Adjusted profit $=$ reported profit $+$ transitory charges $\times(1-t)$ $-$ transitory gains $\times(1-t)$, with $t$ the tax rate applicable to each item.

\qffield{Conditions} IAS~1 is cited as adopted in 2023; for periods from 1~January 2027 IFRS~18 supersedes it. IFRS~18 also requires disclosure of management-defined performance measures that exclude such items (\hyperref[meth:accounting-change-effects]{IFRS~18}). Whether an item recurs is a judgement for the analyst, not a label in the accounts.

\qffield{Practice} \emph{Original} (\hyperref[ex:harbour-tools-statements]{Harbour}, 2025). Adding back the 18 restructuring charge, tax-effected at 25~percent, gives adjusted profit of 68.88 rather than 55.38, and return on equity of 11.00 rather than 8.84~percent. The adjustment is defensible only if restructuring does not recur. A company that restructures every year has a recurring cost whatever its label.

Recasting also captures changes:
\begin{itemize}
\item extending a machine's life from 10 to 14 years after four years spreads the remaining 60 over 10 years, cutting depreciation from 10 to 6 with no restatement;
\item changing the inventory cost formula is a policy change: a 4 lower opening inventory restates opening retained earnings by $-3$ after tax.
\end{itemize}

\qffield{Pitfalls} Adding back ``non-recurring'' charges that recur (big-bath and serial restructurers). Removing charges but keeping the related gains. Forgetting the tax effect. Comparing a restated comparative year with an earlier unrestated year.

\qffield{See also} \hyperref[meth:accounting-change-effects]{effects of accounting changes}; \hyperref[meth:earnings-quality-persistence]{earnings persistence}; \hyperref[meth:earnings-management-presentation]{non-GAAP measures}.

\end{qfmethod}

\begin{qfmethod}[Earnings per share: basic and diluted EPS for simple and complex capital structures, the treasury-stock and if-converted methods, ordering and antidilution]\label{meth:earnings-per-share}
\qffield{Statement} \emph{Basic EPS} is profit attributable to ordinary equity holders of the parent divided by the weighted average number of ordinary shares outstanding \citep[IAS~33 para.~10]{eu-2023-1803-ias33}. The numerator deducts after-tax preference dividends. For cumulative preference shares the deduction is the dividend required for the period, whether declared or not (paras.~12 and~14). Shares are time-weighted by days outstanding (paras.~19--20). Bonus issues and share splits are applied retrospectively to all periods presented (paras.~26 and~64).

\emph{Diluted EPS} adjusts both terms for all dilutive potential ordinary shares (para.~31), deemed converted at the start of the period or at issue if later (para.~36):
\begin{itemize}
\item \emph{options and warrants}, by the treasury-stock method: proceeds are deemed used to buy shares at the period's average market price, and only the shortfall counts as issued for no consideration (paras.~45--46); they dilute only when in the money (para.~47). For share-based payments, the ``price'' includes the fair value of services still to be received (para.~47A);
\item \emph{convertibles}, by the if-converted method: add back after-tax interest or dividends and add the shares on conversion (paras.~33 and~49).
\end{itemize}
Potential shares are antidilutive if they would raise EPS from continuing operations, which is the control number (paras.~41--43). Each issue is tested separately, from the lowest earnings per incremental share to the highest, to maximise dilution (para.~44). Disclosure reconciles numerators and denominators and lists antidilutive instruments (para.~70).

\emph{US GAAP} uses the same structure, presented with equal prominence \citep[p.~374]{financial-statement-analysis-2009}. Since 2020 convertible instruments must use the if-converted method; the treasury-stock method is no longer allowed for them, and share settlement is presumed where cash or shares may be used \citep[pp.~5--6]{us-convertibles-eps-2020}. This applies to SEC-filer public business entities (excluding smaller reporting companies) from fiscal years beginning after 15~December 2021 (p.~6).

\qffield{Notation} Treasury-stock increment $n(1-K/\bar P)$ for $n$ options with exercise price $K$ below the average price $\bar P$. Earnings per incremental share $=\Delta$numerator$/\Delta$shares.

\qffield{Conditions} IAS~33 as adopted in 2023. Control number: continuing operations. The ordering rule matters only when some potential shares are antidilutive in combination.

\qffield{Practice} \emph{Original} (Kestrel, illustrative):
\begin{itemize}
\item Shares were 40m on 1~January, 6m were issued on 1~April and 2m bought back on 1~October; a 1-for-10 bonus issue followed on 30~November. The weighted average is $(40+4.5-0.5)\times1.1=48.4$m, which a month-by-month average confirms.
\item Profit is 50, less a cumulative preference dividend of 2 and a convertible preference dividend of 3, leaving 45. Basic EPS is 0.9298.
\item Vested options over 3m shares at 8, with an average price of 12, add 1m shares (EPI 0). A convertible bond adds 4.5 after-tax interest and 5m shares (EPI 0.90). Including both gives 0.9099.
\item The convertible preference shares (EPI $3/2.5=1.20$) would raise EPS and are antidilutive.
\item Unvested options at 9 with 3.50 of service still to be rendered have an effective price of 12.50, above the average of 12, so they add nothing.
\item A brute-force search over all combinations confirms 0.9099 as the most dilutive figure.
\end{itemize}
\emph{Harbour} 2025: weighted shares $100+0.75\times0.5=100.375$m give basic EPS of 0.5517. Five million vested options at 8, with an average price of 11, add 1.364m shares, giving diluted EPS of 0.5443.

\qffield{Pitfalls} Using year-end rather than weighted shares. Forgetting to restate prior-period EPS for a bonus issue. Including in-the-money options without the treasury-stock offset. Treating EPS growth from buybacks as operating improvement (\hyperref[meth:return-on-equity-decomposition]{ROE decomposition}).

\qffield{See also} \hyperref[meth:share-based-compensation-modelling]{forecasting dilution}; \hyperref[meth:income-statement-ratios]{income-statement ratios}.

\end{qfmethod}

\begin{qfmethod}[Common-size income statements and income-statement ratios: margins, expense structure, operating leverage, and evaluating performance across years and companies]\label{meth:income-statement-ratios}
\qffield{Statement} A common-size income statement expresses every line as a percentage of revenue. It shows the cost structure and how it changes, and it allows companies of different size to be compared \citep[pp.~28--33]{financial-statement-analysis-2009}. Year-to-year change analysis, in amounts and percentages, complements it. Percentage changes from negative, zero or very small bases are not meaningful (p.~28). Margins, such as gross, operating and net profit over revenue, are interpreted against the company's history, predetermined standards and competitors, and their variability matters as well as their trend (p.~35).

Classification affects comparability. IAS~1 lets entities analyse expenses by nature or by function, the latter needing allocation judgements \citep[IAS~1 paras.~99 and~103]{eu-2023-1803-ias1}. From 2027 IFRS~18 requires operating profit and profit before financing and income taxes as subtotals \citep[Annex, IFRS~18 para.~69]{eu-2026-338-ifrs18}. A function presentation must then present cost of sales separately and disclose, in one note, total depreciation, amortisation, employee benefits, impairments and inventory write-downs (paras.~82--83). US public business entities will disclose purchases of inventory, employee compensation, depreciation and intangible amortisation within each relevant expense caption, plus total selling expenses, for annual periods beginning after 15~December 2026 \citep[pp.~1--3]{us-expense-disaggregation-2024}.

\qffield{Notation} Gross margin $=(\text{revenue}-\text{cost of sales})/\text{revenue}$; operating margin $=$ operating profit$/$revenue; net margin $=$ profit$/$revenue; effective tax rate $=$ tax$/$profit before tax.

\qffield{Conditions} Margins are only comparable when classification is: one company's cost of sales may hold expenses another reports in administration. Different revenue recognition (\hyperref[meth:revenue-recognition-analysis]{gross versus net}) changes every margin.

\qffield{Practice} \emph{Original} (\hyperref[ex:harbour-tools-statements]{Harbour}).

\begin{center}\small
\begin{tabular}{@{}lrr@{}}
\toprule
Percent of revenue & 2025 & 2024\\
\midrule
Cost of sales & 64.54 & 64.00\\
Gross profit & 35.46 & 36.00\\
Selling and administration & 21.57 & 21.80\\
Research and development & 3.52 & 3.50\\
Restructuring & 1.67 & 0.00\\
Operating profit & 8.61 & 11.00\\
Finance costs & 1.50 & 1.80\\
Income tax & 1.99 & 2.50\\
Profit & 5.13 & 6.70\\
\bottomrule
\end{tabular}
\end{center}

\emph{Evaluation.} Revenue rose 8~percent, but operating margin fell 2.39 points:
\begin{itemize}
\item 1.67 points came from restructuring;
\item 0.54 points came from gross margin, as cost of sales grew 8.9~percent against revenue of 8~percent, so price increases did not fully offset cost inflation;
\item 0.39 points came from the swing from a disposal gain of 3 in 2024 to a loss of 1;
\item selling and administrative costs fell by 0.23 points as a share of revenue, a modest operating-leverage benefit.
\end{itemize}
Net margin fell less, by 1.57 points, because finance costs and tax were lower relative to revenue.
The effective tax rate rose from 27.17 to 27.93~percent because the non-deductible share-based expense of 9 weighs more on a smaller profit. Excluding restructuring, operating margin was 10.28~percent, still below 2024. The gross-margin pressure is the issue that the forecast (\hyperref[meth:competition-inflation-forecasting]{pass-through}) must address.

\qffield{Pitfalls} Reading a lower expense ratio as efficiency when it reflects cuts to discretionary spending such as advertising or research, which harm future sales \citep[p.~35]{financial-statement-analysis-2009}. Comparing margins across companies that classify expenses differently. Ignoring the tax-rate effect of permanent differences.

\qffield{See also} \hyperref[meth:financial-ratio-families]{ratio families}; \hyperref[meth:return-on-equity-decomposition]{ROE decomposition}; \hyperref[meth:tax-rate-measures]{tax-rate measures}.

\end{qfmethod}
